\answer
\begin{enumerate}

%%--- 章末問題 2.1 ---
\item[\ref{ex2.1}]
(a) 多数キャリアの電子は$n \approx N_d = 10^{17}\,\mathrm{cm^{-3}}$。
少数キャリアの正孔は質量作用の法則\ref{eq2.1}より
\begin{equation*}
p = \frac{n_i^2}{n} = \frac{(1.5\times10^{10})^2}{10^{17}}
\approx 2.3\times10^{3}\,\mathrm{cm^{-3}}
\end{equation*}

(b) 式\ref{eq2.2}で正孔の項を無視して
\begin{equation*}
\sigma = e n \mu_n = 1.602\times10^{-19}\times10^{17}\times1350
\approx 22\,\mathrm{S/cm}\ (= 2.2\times10^{3}\,\mathrm{S/m})
\end{equation*}

%%--- 章末問題 2.2 ---
\item[\ref{ex2.2}]
接合面ではキャリアの濃度が急変しているので，
n形側の電子はp形側へ，p形側の正孔はn形側へ\textbf{拡散}する。
流れ込んだ先は相手のキャリアだらけだから，電子と正孔は出会って対で消滅する（\textbf{再結合}）。
キャリアが消えた接合面付近には，n形側にドナーイオン（$+$），p形側にアクセプタイオン（$-$）が残る。
イオンは結晶に組み込まれていて動けない\textbf{固定電荷}であり，
正負の固定電荷が向かい合うことでn形からp形へ向かう\textbf{内蔵電場}ができる。
内蔵電場は電子をn形側へ，正孔をp形側へ押し戻す向き，つまり拡散をせき止める向きに働く。
拡散が進むほど固定電荷が増えて内蔵電場が強まり，
やがて拡散と電場による引き戻しがつり合って，
動けるキャリアのほとんどない層，すなわち空乏層が一定の幅で残る。

%%--- 章末問題 2.3 ---
\item[\ref{ex2.3}]
(a) 真性キャリア濃度の観点：
バンドギャップが大きいほど，熱で価電子帯から伝導帯へ上がる電子は
急激に少なくなる（\ref{sec2.1}節）。
真性キャリア濃度が小さいので，オフ状態の漏れ電流が小さく，
高温になってもドーピングで作ったキャリア数が熱で生まれるキャリアに埋もれない。
つまり高温環境（自動車・産業機器）でも確実にオフを保てる。

(b) 絶縁破壊電界の観点：
バンドギャップが大きい材料は絶縁破壊電界も大きい
（SiCはシリコンの約10倍）。
同じ耐圧を得るのに必要な空乏層（ドリフト層）を薄く，
かつ高濃度にできるため，例題\ref{rei2.2}で見たドリフト層の抵抗，
すなわちオン抵抗を大幅に下げられる。
耐圧とオン抵抗のトレードオフそのものが材料の性質で緩和される。
これがSiC・GaNが「高耐圧と低損失の両立」を実現できる理由である（3章）。

\end{enumerate}
